
\begin{abstract}

Simultaneous speech translation (SST) produces target text incrementally from partial speech input. 
Recent speech large language models (Speech LLMs) have substantially improved SST quality, yet they still struggle to correctly translate rare and domain-specific terminology. 
While retrieval augmentation has been effective for terminology translation in machine translation, bringing retrieval to SST is non-trivial: it requires fast and accurate cross-modal (speech-to-text) retrieval under partial, continually arriving input, and the model must decide whether and when to apply retrieved terms during incremental generation.
We propose \textbf{R}etrieval-\textbf{A}ugmented \textbf{S}imultaneous \textbf{S}peech \textbf{T}ranslation (\method), which tightly integrates cross-modal retrieval into the SST pipeline. 
\method trains a lightweight speech--text retriever and performs efficient sliding-window retrieval, providing chunkwise terminology hints to the Speech LLM. 
We further synthesize training data that teaches the Speech LLM to leverage retrieved terms precisely. 
Experiments on three language directions of the ACL 60/60 dev set show that \method improves terminology translation accuracy by up to \textbf{16\%} and increases overall translation quality by up to \textbf{3 BLEU} points, with ablations confirming the contribution of each component.

\end{abstract}